\documentclass[11pt,a4paper]{article}
\usepackage{amsmath,amssymb,amsthm,graphicx}
\newtheorem{theorem}{Theorem}
\theoremstyle{definition}
\newtheorem{definition}[theorem]{Definition}
\newcommand{\norm}[1]{\left\lVert #1\right\rVert}
\newcommand{\inner}[2]{\langle #1,#2\rangle}

\title{Orthogonal Projections in Hilbert Space}
\author{Elena Marchetti}
\date{February 27, 2024}

\begin{document}
\maketitle
\tableofcontents

\section{Preliminaries}
Let $H$ be a real Hilbert space with inner product $\inner{\cdot}{\cdot}$ and norm $\norm{x}=\sqrt{\inner{x}{x}}$. The \textsc{Cauchy--Schwarz} inequality $|\inner{x}{y}|\le\norm{x}\norm{y}$ holds for all $x,y\in H$.

\begin{definition}
A subset $C\subseteq H$ is \emph{convex} if $tx+(1-t)y\in C$ for all $x,y\in C$ and $t\in[0,1]$.
\end{definition}

\section{The projection theorem}
\begin{theorem}\label{thm:proj}
Let $C\subseteq H$ be nonempty, closed and convex, and let $x\in H$. Then there is a unique $p\in C$ with $\norm{x-p}=\min_{y\in C}\norm{x-y}$, and $p$ is characterised by
\begin{equation}\label{eq:var}
  \inner{x-p}{y-p}\le0\quad\text{for all } y\in C.
\end{equation}
\end{theorem}

\begin{proof}
Let $d=\inf_{y\in C}\norm{x-y}$ and choose $y_n\in C$ with $\norm{x-y_n}\to d$. The parallelogram law gives
\begin{multline*}
  \norm{y_n-y_m}^2 = 2\norm{x-y_n}^2+2\norm{x-y_m}^2\\
  - 4\Bigl\lVert x-\frac{y_n+y_m}{2}\Bigr\rVert^2 \le 2\norm{x-y_n}^2+2\norm{x-y_m}^2-4d^2 ,
\end{multline*}
which tends to $0$, so $(y_n)$ is Cauchy and converges to some $p\in C$. For \eqref{eq:var}, apply minimality to $p+t(y-p)\in C$ and let $t\downarrow0$.
\end{proof}

\subsection{The linear case}
If $C=M$ is a closed subspace, \eqref{eq:var} forces $\inner{x-p}{y}=0$ for all $y\in M$, so $x-p\in M^\perp$ and $H=M\oplus M^\perp$.

\subsubsection{Example}
In $\mathbb{R}^3$ take $M=\operatorname{span}\{(1,1,0),(0,1,1)\}$. The projection of $x$ onto $M$ is $P x$, where $P=A(A^{\mathsf T}A)^{-1}A^{\mathsf T}$ and the columns of $A$ span $M$.

\begin{figure}[b]
\centering
\includegraphics[width=0.25\textwidth]{images/square.png}
\caption{Sketch of a point and its projection onto a plane.}
\label{fig:proj}
\end{figure}

See Figure~\ref{fig:proj} on page~\pageref{fig:proj}.

\begin{thebibliography}{9}
\bibitem{rudin} W.~Rudin, \emph{Real and Complex Analysis}, McGraw-Hill, 1987.
\bibitem{brezis} H.~Brezis, \emph{Functional Analysis, Sobolev Spaces and Partial Differential Equations}, Springer, 2011.
\end{thebibliography}

The proofs follow the presentation in \cite{brezis}; the reader may also consult \cite{rudin}.

\end{document}
